% ИСТОЧНИК (как прислан владельцу, дословно). Листок собеседования в 8 класс, школа 179.
% Формат: 3 часа, дети после 7 класса; на руках — статья «Кривые дракона» («Квант» 2020 №1).
% Даже за 3 часа группа не прошла листок целиком. Наш блок: 4 занятия × 1,5 ч = 6 ч, дети после 6 класса.
% Закомментированные (%) куски — так было в присланном исходнике: часть упражнений статьи отключена,
% часть задач помечена как «прочитать, но не сдавать». Комментарии внутри — авторские, НЕ наши.
% NB: счётчики \setcounter{z}{...} задают нумерацию, поэтому видимых задач меньше, чем \z в файле.

\centerline{\textsc{\Large Задачи про ломаные дракона}}

Упражнения из статьи № 3а, 4 и 5б являются задачами,  их можно сдавать. Остальные упражнения  советуем прочитать. Вы можете обсудить их с принимающим, но плюсики за них не ставятся.
Можно пользоваться теоремами из статьи (даже если вы не до конца понимаете их доказательства).

% \z  Какой длины надо взять полоску, чтобы, сложив ее пополам 30 раз, получить расстояние между соседними сгибами равным 1 см? Больше или меньше расстояния от Земли до Луны?


% \z  Как изменится ломаная дракона, если полоску бумаги положить на стол другим ребром? Как изменится при этом слово, записывающее ломаную?


 % \z \leth [Мне кажется, эта задача нужна для того, чтобы разобраться в тексте, мы хотим её проверять? или оставить как совет?] Сколько существует различных (не подобных друг другу) ломаных дракона ранга 4? Нарисуйте их все. Напишите соответствующие им слова из букв $L$ и $R$.
% \lett  Сколько существует различных ломаных ранга $n$?


% \z Допустим, что черепаха проползла по ломаной дракона и прочла слово из букв $L$ и $R$. Какое слово она прочтет, если проползет по этой ломаной в обратном направлении?

% \z  Пусть $s$~--- некоторое слово из букв $L$ и $R$. Обозначим через $\overline{s}$ слово, которое получится из $s$, если переставить в нем буквы в обратном порядке и потом поменять $L$ на $R$ и $R$ на $L$. (Например, если $s = LLR$, то $\overline{s} = LRR$.)
% Пусть $s$ и $t$~--- некоторые слова из букв $L$ и $R$. Будем обозначать через $st$ слово, получающееся, если слова $s$ и $t$ написать рядом.
% Докажите, что:
% \lett  Для любых слов $s$ и $t$ $$\overline{st}=\overline{t}\overline{s}.$$
\setcounter{z}{4}
\z \setcounter{lett}{1}
% \leth Пусть $s_{11}$~--- слово, соответствующее произвольной ломаной дракона ранга $11$. Объясните, как составить это слово $s_{11}$ из некоторого слова $s_{10}$, соответствующего ломаной дракона ранга ранга 10, сопряжённого ему слова $\overline{s_{10}}$ и, возможно, дополнительных букв $L$ и $R$ или слов другого ранга.
\leth где $s_n$~---  слово, соответствующее какой-то ломаной дракона ранга~$n$, а $x$~--- буква $L$ или $R$.
\lett  Пусть $s$~--- слово, соответствующее ломаной дракона. Чем отличаются между собой слова  $s$ и $\overline{s}$?
\lett Известно, что в слове, записывающем ломаную дракона
ранга 8, на 1-м, 2-м, 4-м, 8-м, 16-м, 32-м, 64-м и  128-м  местах стоят соответственно буквы $L, L, R, R, R, L, R, L$. Какая буква стоит на
179-м месте?
\leth А на 180-м?
% \lettstar Можно ли определить все буквы слова из  пункта (г)? Если да, то как?


\z Из ломаной дракона ранга $n$ можно получить две ломаные дракона ранга $n+1$, пользуясь либо теоремой 1, либо теоремой 2.
\lett Будут ли в обоих случаях всегда получаться те же самые две ломаные дракона ранга $n+1$ или, вообще говоря, другие?
\lett [\textbf{Письменно}] Как получать их слова из слова данной ломаной ранга $n$?

\z Может ли в слове, соответствующем ломаной дракона, быть 4 одинаковых буквы подряд?

\z Найдите  середину ломаной дракона, изображённой на обороте, а именно:
\lett объясните, как построить точку, в которой находится середина;
\lett укажите конкретный изгиб ломаной, в котором она находится;
\lett Какого ранга эта ломаная дракона?
% Эта кривая дракона имеет специальное название <<Мама, папа и сын>>.




\z Введём на клетчатой бумаге систему координат, направив оси по линиям сетки и взяв за единицу масштаба сторону клетки. Будем рисовать Главную ломаную дракона так, чтобы первыми тремя ее вершинами были точки $(0;0)$, $(0;1)$, $(1;1)$. Пользуясь теоремой 1, можно продолжать эту ломаную сколько угодно раз: мы будем получать последовательно ломаные ранга 1, 2, 3,\ldots %Можно представить себе, что мы нарисовали их все и получили бесконечную ломаную дракона. %(главную ломаную дракона ранга $\infty$).
%Её первые $2^n$ звеньев образуют ломаную ранга $n$.
 Найдите координаты конца такой ломаной ранга:
 \leth 5; \leth 12; \leth 179.

\z На плоскости даны  точки $A$ и $B$. Начертим квадрат с центром в середине отрезка $AB$, одна из сторон которого равна $2AB$ и параллельна  $AB$.
\lett Докажите, что любая ломаная дракона с концами в точках $A$ и $B$ лежит внутри этого квадрата.
\lett Верно ли, что для любого квадратика, содержащегося в нашем большом квадрате, найдётся ломаная дракона с концами $A$ и $B$, <<залезающая>> в этот квадратик (то есть проходящая хотя бы через одну его точку)?

% Докажите, что для любой точки М внутри этого квадрата и любого положительного числа  можно найти такую ломаную дракона с концами А и В, которая проходит от точки М на расстоянии меньше  (говоря математическим языком, «объединение всех ломаных дракона с конца- ми в точках А и В всюду плотно в этом квадрате»).

% \z Пусть черепаха ползет по плоскости из точки $A$ с постоянной скоростью --- вначале по направлению $AB$, а затем через каждые 15 минут поворачивает на $90^{\circ}$. Докажите, что вернуться в $A$ черепаха может лишь через целое число часов и лишь по направлению, перпендикулярному к $AB$.


\z  Придумайте  про ломаные дракона ваши собственные:
\lett вопрос; \leth гипотезу; \leth теорему.
% \lett Придумайте  гипотезу про ломаные дракона.
% \lett Докажите свой факт про ломаную дракона (то есть придумайте свою теорему)


\newpage

\

\vfill

\begin{center}
    \includegraphics[width=.9\linewidth]{Figs/dragon_family2.png}
\end{center}


\vfill

\

% \lett Слово, являющееся записью этой ломаной, можно быстро написать так. Сначала записываем чередующиеся буквы $L$ и $R$, оставляя между ними промежутки для пропущенных букв: $L\square R\square L\square R\square L\square R\square L\square R\ldots$, а затем поступаем следующим образом. Пальцем левой руки укажем на первую букву, а правой рукой запишем ее в первый промежуток, затем левой рукой укажем на вторую букву, а правой вставим ее во второй промежуток, левой на третью\ldots и так последовательно по всем буквам (не пропуская вновь вставленных), а правой вносим каждую указанную букву в первый свободный промежуток:
% $LLRLLRRLLLRRLRRLLLRL\ldots$



% \z В каком случае два разных слова задают одинаковые по форме (подобные) кривые дракона? Опишите все такие пары слов.



\newpage

\centerline{\textsc{\Large Ещё задачи}}



\setcounter{z}{9}
\z \setcounter{lett}{2}
\leth Представьте, что мы нарисуем все ломаные дракона с концами в точках $A$ и $B$. Какую фигуру они <<заполняют>>?


\setcounter{z}{11}


\z Раскрасим клетчатую плоскость в шахматном порядке. Теперь на сторонах каждой белой клетки нарисуем стрелочки так, чтобы они обходили её против часовой стрелки.
\lett Докажите, что эти стрелки обходят чёрные клетки по часовой стрелке.
\lett Ломаная дракона нарисована на этой клетчатой плоскости так, что первое её звено совпадает со стороной какой-то клетки, а стрелка на этой стороне ведёт из конца ломаной в направлении её середины. Докажите, что если черепаха будет ползти по ломаной дракона, начиная из этого конца, то она всегда будет ползти в направлении стрелок.

\z Докажите, что ломаная дракона никогда не проходит по одному и тому же отрезку более одного раза.


\z Решите задачу со стр. 19 статьи.


\vfill

% Далее в исходнике — повтор блока «Ещё задачи» (страница-дубль) и страницы с картинкой Figs/паровоз.png (×4).
